\documentclass[11pt]{article}
\usepackage[utf8]{inputenc}
\usepackage{amsmath}
\usepackage{lipsum} 
\usepackage{fullpage}
\usepackage{mathrsfs}
\usepackage{svg}
\usepackage{xr}
\usepackage{amsthm,cite,natbib,url,graphicx,booktabs,lipsum,color,bm,caption,subcaption,soul}
\usepackage{pifont,tikz,paralist,multirow,amssymb}
\usepackage{xifthen}
\usepackage{enumerate}
\usepackage{enumitem}
\usepackage{titlesec}
\usepackage{xcolor}
\definecolor{revisionmodified}{HTML}{F2A6A6}   % red: modified text
\definecolor{revisionadded}{HTML}{8AC1AF}      % green: newly added text
\definecolor{revisionmoved}{HTML}{B4D9FF}      % blue: relocated/moved text

% 自定义 Reviewer 计数器
\newcounter{reviewer}
\setcounter{reviewer}{0}
\newcounter{point}[reviewer]
\setcounter{point}{0}

% % 将 Reviewer 的计数方式改为字母（A, B, C, D, ...）
% \renewcommand{\thereviewer}{\Alph{reviewer}}

% 定义 Reviewer 的引用格式
\renewcommand{\thepoint}{Comment\,\thereviewer.\arabic{point}}

% 定义 Associate Editor 的评论环境
\newenvironment{point1}
   {\refstepcounter{point} \bigskip \noindent {\textbf{Editor Comment~}} ---\ }
   {\par }

% 定义 Reviewer 的评论环境
\newenvironment{point}
   {\refstepcounter{point} \bigskip \noindent {\textbf{Reviewer~\thepoint} } ---\ }
   {\par }

% 定义 Reviewer 的简短评论命令
\newcommand{\shortpoint}[1]{\refstepcounter{point}   \bigskip \noindent 
	{\textbf{Reviewer~Point~\thepoint} } ---~#1  }

% 定义回复环境
\newenvironment{reply}
   {\medskip \noindent \color{blue} \begin{sf}\textbf{Reply}:\  }
   {\medskip \par \end{sf}}

% 定义简短回复命令
\newcommand{\shortreply}[2][]{\medskip \noindent \begin{sf}\textbf{Reply}:\  #2
	\ifthenelse{\equal{#1}{}}{}{ \hfill \footnotesize (#1)}%
	\medskip \end{sf}}

% 允许跳过某些 Reviewer
\newcommand{\skipreviewer}{\stepcounter{reviewer}}

% 定义 Reviewer 部分
\newcommand{\reviewersection}{\stepcounter{reviewer} \bigskip \hrule
                  \section*{Reviewer \thereviewer}}

\title{Response to Reviewers \\ Manuscript ID: CEUS-D-25-01256}
\author{}
\date{}

\begin{document}
\maketitle

\vspace{-30pt}

The authors would like to thank the editor and reviewers for their constructive comments and suggestions, which have greatly improved the quality of this manuscript. In response, we have carefully revised the manuscript in accordance with all editorial and reviewers’ comments. Our point-by-point responses are provided below. In the revised manuscript, all changes made in response to the editor’s and reviewers’ comments are highlighted as follows: 
\colorbox{revisionmodified}{\strut modified content}, 
\colorbox{revisionadded}{\strut newly added content}, and 
\colorbox{revisionmoved}{\strut content relocated within the manuscript}.

\begin{point1}
On top of the suggestions by anonymous reviewers, please consider expanding the data availability statement by including an archived version of the \texttt{SpatLyu/tEDM\_CEUS} and \texttt{stscl/tEDM} repositories, and link the DOI from the statement alongside the GitHub URL. An easy way to do that is to link your repositories to Zenodo and archive released versions there. That way we can ensure longevity of links and point to the exact version of the contents you refer to in the paper. Otherwise, the contents of the repository may change and get out of sync with the contents of the paper.

Please note that submissions to the SI do not need to be anonymised, so feel free to skip spending time anonymising the revision.
\end{point1}

\begin{reply}

\end{reply}

\paragraph{Note:} All line numbers mentioned in this response letter refer to the small line numbers adjacent to the text in the revised manuscript (as generated by the \LaTeX{} line-numbering package). These may not correspond to the line numbers used by the submission system, which tend to be larger and differ from those in the manuscript.

\reviewersection

\begin{point}
\textbf{Novelty and Methodological Contribution}

The proposed tool appears to mainly integrate several existing EDM methods without introducing methodological improvements to address known limitations of these models. If the contribution lies primarily in integration, the level of innovation may be limited. In particular, the authors claim that tEDM extends EDM to multivariate and spatially replicated time series, but the manuscript does not describe how this extension is implemented or evaluated in detail. More clarification on the underlying methodological framework is necessary.
\end{point}

\begin{reply}
We thank the reviewer for this constructive comment. We would like to clarify that \textit{tEDM} is not a simple integration of existing EDM methods, but introduces several methodological extensions to address known limitations of CCM, PCM, and CMC in the context of urban time series. In the revised manuscript, we expanded the background and methodological descriptions to explicitly articulate these innovations and their implementation details (see revised manuscript, lines~xx--xx and xx--xx).

First, building upon multispatial CCM and related advances, \textit{tEDM} extends CCM to explicitly leverage spatially replicated observations that frequently arise in urban monitoring networks. Replicated time series are independently embedded, temporally aligned, and concatenated to form a combined embedding, enabling causal detection when individual series are short or weakly informative. Spatial dependence among replicated units can be incorporated through informed selection of spatial neighborhoods, improving effective sample size while preserving dynamical consistency.

Second, PCM is generalized beyond separating indirect effects to focus on identifying direct causal influences in high-dimensional systems. In \textit{tEDM}, potential third-party variables, including mediator, collider, confounding structures, and their combinations, are treated as conditioning variables, allowing the algorithm to isolate direct causal pathways under complex multivariate dependencies.

Third, for causal strength estimation, CMC is augmented with a statistically robust DeLong placement strategy, which enables consistent estimation of causal strength together with uncertainty quantification via confidence intervals and significance testing. This extension improves the reliability of causal inference under noisy urban data conditions.

These methodological extensions are described both conceptually and algorithmically in the revised manuscript, and their effectiveness is demonstrated through the three case studies covering different urban domains. We believe these additions address the reviewer’s concern regarding methodological novelty and clarify how the proposed extensions are implemented and evaluated.
\end{reply}

\begin{point}
\textbf{Adaptation to Urban Environmental Research}

The manuscript repeatedly states that tEDM is specifically developed for urban environmental studies. However, there is no clear explanation of how the model has been adapted to the specific characteristics of urban data, nor how spatial dependencies are incorporated into the modeling process. The authors should elaborate on why this tool is particularly suited for urban environments and provide supporting analysis or discussion to justify this claim.
\end{point}

\begin{reply}

\end{reply}

\begin{point}
\textbf{Comparison with Other Causal Discovery Methods}

The paper does not clearly discuss how the causal discovery results produced by tEDM differ from those obtained using other widely adopted causal discovery algorithms, such as the PC algorithm. It would be valuable for readers if the authors could discuss which approach is more suitable under certain conditions, and—given that the true causal structure is typically unknown—how the reliability and robustness of the inferred causal networks are evaluated.
\end{point}

\begin{reply}

\end{reply}

\begin{point}
\textbf{Minor Comments on Presentation}

\begin{enumerate}[label=\arabic*.]
    \item In Figure 4, it is not clearly explained how the causal network results shown in Figure 3-c were obtained. This connection is crucial for understanding the workflow and logic of the causal discovery process. Please provide additional clarification or methodological description.
    \item In Figure 3-b, the meaning of the symbol ``\#'' is unclear. A brief note or legend explanation should be added to help readers interpret the figure correctly.
\end{enumerate}
\end{point}

\begin{reply}

\end{reply}

\reviewersection 

\begin{point}
The authors need to further explain why the results obtained from the tEDM analysis represent temporal causality rather than merely correlation.
\end{point}

\begin{reply}
We thank the reviewer for raising this important conceptual clarification. In the revised manuscript, we substantially expanded the \emph{Discussion and conclusion} section to explicitly explain why the causal relationships inferred by \textit{tEDM} represent \emph{temporal and dynamical causality} rather than mere statistical correlation (see revised manuscript, lines~xx--xx).

Specifically, we now emphasize that \textit{tEDM} is grounded in the concept of dynamical causality, where causation is defined by whether the evolution of one variable leaves a persistent and directionally identifiable imprint on the state evolution of another variable through the underlying system dynamics. Unlike correlation analysis, which measures instantaneous or marginal association, EDM-based methods exploit state-space reconstruction and cross mapping to test whether information about a potential driver can be reliably recovered from the temporal trajectory of a response variable. Moreover, the convergence of cross-mapping skill with increasing library size provides an operational criterion to distinguish genuine causal signal from spurious correlation induced by shared trends, synchronization, or external forcing.
\end{reply}

\begin{point}
Complex urban systems involve variables across multiple dimensions such as population, economy, transportation, energy, and environment. These variables inherently interact in a disordered, complex, and stochastic manner. How does tEDM characterize the features of different variables and adapt to various scenarios?
\end{point}

\begin{reply}
We appreciate the reviewer's insightful question regarding the adaptability of \textit{tEDM} to heterogeneous and high-dimensional urban systems. In the revised manuscript, we clarified that \textit{tEDM} does not rely on predefined parametric assumptions or variable-specific feature engineering. Instead, it characterizes variables through their observed dynamical behavior using state-space reconstruction, allowing heterogeneous variables (e.g., environmental, socioeconomic, epidemiological, and infrastructural indicators) to be analyzed within a unified dynamical framework (see revised manuscript, lines~xx--xx).

Adaptability across scenarios is achieved through several mechanisms. Spatially replicated observations allow the method to exploit shared dynamical structure across locations, alleviating data sparsity and heterogeneity in short time series. Multivariate conditioning in PCM enables the algorithm to disentangle direct causal influences from indirect, confounded, or collider-driven dependencies in complex interaction networks. The geometric criterion in CMC provides robustness when prediction-based metrics become unstable under noise or partial observability. Together, these complementary strategies allow \textit{tEDM} to adapt flexibly to diverse data regimes, noise levels, dimensionalities, and coupling structures commonly encountered in urban systems.

We further emphasize that the data-driven and nonparametric nature of EDM enables \textit{tEDM} to accommodate nonlinear, nonstationary, and strongly coupled dynamics without requiring explicit model specification or variable-specific assumptions. The presented case studies across environmental health, climate–energy interaction, and epidemic spreading demonstrate the practical adaptability of the framework to heterogeneous urban scenarios.
\end{reply}

\begin{point}
As an R package, the authors should present the basic framework of the package as well as the data flow.
\end{point}

\begin{reply}
We thank the reviewer for highlighting the need to clarify the overall framework of the package and the organization of the software. In the revised manuscript (Lines~xx--xx), we have substantially expanded the subsection 3.1 ``Software design philosophy and architecture'' to explicitly describe the layered architecture of \textit{tEDM}, including the separation between the user-facing R interface layer and the high-performance C++ computational backend. We further added a new schematic figure (Figure~\ref{fig_tedm_arch}) illustrating the source code organization and architectural layers of the package, which provides a clear overview of the basic framework and module structure.

The description now explains the responsibilities of each layer, including user interaction, data preparation, parameter configuration, result formatting, visualization at the R level, and computationally intensive numerical operations at the C++ level. To avoid redundancy and improve clarity, the detailed data flow and execution pipeline are presented separately in the following section of the revised manuscript (Lines~xx--xx).
\end{reply}

\begin{point}
Which existing R packages does this package rely on, and how does it ensure computational efficiency when handling large-scale data?
\end{point}

\begin{reply}
We appreciate the reviewer's suggestion to clarify the software dependencies and the mechanisms ensuring computational efficiency for large-scale data analysis. In the revised manuscript (Lines~xx--xx), we explicitly document the mandatory software stack of \textit{tEDM}. On the computational side, the package relies on \textit{Rcpp}, \textit{RcppArmadillo}, and \textit{RcppThread} to enable efficient C++ integration, optimized linear algebra routines, and multithreaded parallel execution. Core algorithmic components such as state-space reconstruction, nearest-neighbor search, and cross mapping are implemented in C++ using the Standard Template Library (STL), thereby minimizing runtime overhead and improving scalability.

On the user interface side, \textit{ggplot2} is employed for visualization and \textit{dplyr} is used for data manipulation and preprocessing, ensuring seamless interoperability with existing R-based workflows. The revised text also clarifies that all computationally intensive operations are executed in the optimized C++ backend, while the R layer primarily manages data input, configuration, and result presentation. These additions directly address how \textit{tEDM} achieves both usability and computational efficiency when handling large-scale urban time series data.
\end{reply}

\begin{point}
In Figure 5, what is the meaning of ``causal strength''? Does ``causal strength'' have general applicability?
\end{point}

\begin{reply}
We appreciate the reviewer's request for clarification regarding the meaning and applicability of ``causal strength.'' In the revised manuscript, we added a dedicated paragraph in the \emph{Discussion and conclusion} section to formally define this concept and clarify its scope of interpretation (see revised manuscript, lines~xx--xx).

We now explain that within the dynamical causality framework, quantities reported as ``cross-mapping skill'' (the $\rho$ value) in CCM and PCM, and ``causal strength'' in CMC quantify the influence of data-driven deterministic processes governing the evolution of the underlying system. These metrics measure the consistency and reliability with which the state of one variable can be recovered from another in reconstructed state space, thereby serving as relative indicators of directional dynamical dependence.

Importantly, we clarify that ``causal strength'' should not be interpreted as a counterfactual causal effect comparable to quantities such as the average treatment effect in intervention-based causal inference. Its magnitude is system-dependent and is primarily intended for causal discovery purposes, including ranking, comparison, and screening of potential causal influences within a given dataset, rather than direct transferability across unrelated systems or policy contexts. We further note that in applied analyses, statistical significance should be accompanied by effect-size thresholds (e.g., $\rho > 0.2$ in our case studies) to reduce spurious detections caused by strong coupling or mirror dynamics.
\end{reply}

\baselineskip12pt
\bibliographystyle{elsarticle-harv} 
\bibliography{references.bib} 

\end{document}